Among various bandit algorithms, there are two most commonly adopted online bandit strategies. One explores based on the reward estimates sampled from context-action pairs' posterior distributions, represented by Thompson sampling \cite{thompson1933likelihood}, and the other follows the "optimism in the face of uncertainty" principle, represented by upper confidence bound (UCB) \cite{agrawal1995sample, auer2002finite}. In this section, we offer a high level introduction to Thompson sampling and UCB algorithms as well as their extensions to contextual settings. The deep neural network versions of the algorithms are considered baselines for our work. We will also review related industry and practical adoptions of these strategies. Note that in this section, we limit our discussion to exploration strategy for immediate reward and do not cover optimization approaches for cumulative returns \cite{zhu2023deep, xu2023optimize, chen2021exploration}.

\subsection{Thompson Sampling, Its Extensions and Related Adoptions} \label{sec:ts}
Thompson sampling is an exploration strategy that samples reward estimates from context-action pairs' posterior distributions and then executes greedily with respect to these samples. This strategy is particularly popular due to its simplicity. Following the definition in Section \ref{sec:prob}, we first consider a vanilla Thompson sampling algorithm \cite{thompson1933likelihood} for multi-armed bandits. For a multi-armed bandit, the context $S_t$ is always the same and the agent needs to choose $A_t$ from $\mathcal{A}$ without any representation learning. A Thompson sampling agent keeps a reward posterior distribution $\hat{P}_{a}$ for each arm $a$ and updates it over time. At time step $t$, the agent chooses
\begin{equation}
    A_t = \arg\max_{a \in \mathcal{A}} \hat{R}_{t, a}, \qquad \hat{R}_{t, a} \sim \hat{P}_a,
\end{equation}
where $\hat{R}_{t, a}$ is the estimated reward sampled from the agent's posterior belief. The agent then updates its posterior belief $\hat{P}_{a}$ with the latest reward $R_{t+1}$ and selected action $A_t$. The vanilla Thompson sampling approach is a popular production strategy when it comes to small action space recommendations \cite{sanz2019simple, hsieh2015efficient}. 

Extending Thompson sampling to large action space and contextual bandits, an agent can compute parametric estimates of rewards of context-action pairs by sampling linear model parameters from the posterior belief \cite{abeille2017linear, agrawal2013thompson,russo2013eluder,russo2014posterior,russo2018tutorial} instead of sampling a point estimate from the reward posterior distribution. Here, a linear Thompson sampling agent keeps track of a posterior belief $\hat{P}_\theta$ over the parameters $\theta$. The agent chooses action
\begin{equation}
    A_t = \arg\max_{a \in \mathcal{A}} \hat{\theta}_t^T x_{t, a}, \qquad \hat{\theta}_t \sim \hat{P}_\theta, x_{t, a} = \text{concat}(\psi_{S_t}, \phi_{a}). 
\end{equation}
The agent then updates its posterior belief of $\hat{P}_\theta$ with the latest reward $R_{t+1}$, context $S_t$ and action $A_t$. The posterior update could be performed via Laplace approximation \cite{laplace1986memoir}, and is also presented in industry adoptions of linear Thompson sampling \cite{chapelle2011empirical}. 

The method above could be extended to neural networks as the agent can also keep a posterior belief over neural network parameters. To avoid heavy computation of Laplace approximation, various Bayesian methods have been proposed to achieve posterior updates for distributions over neural network parameters, including deep ensemble \cite{lakshminarayanan2017simple}, ensemble sampling with prior networks \cite{lu2017ensemble}, Bayes by Backprop \cite{blundell2015weight}, Hypermodels \cite{dwaracherla2019hypermodels}, Monte-Carlo Dropouts \cite{gal2016dropout} and EpiNet \cite{osband2021epistemic}. Among the methods above, the two ensemble methods \cite{broden2018ensemble, lu2018efficient, tang2014ensemble} and Monte-Carlo Dropout \cite{guo2020deep} are the most commonly adopted methods for contextual bandits with parameter posterior sampling in a practical setting.

One additional line of work that leverages Thompson sampling for neural networks is neural Thompson sampling \cite{zhang2021neural}. In this approach, instead of sampling from parameters' posterior distribution, neural Thompson sampling directly samples from the posterior estimate of the reward of a context-action pair. A Neural Thompson sampling agent keeps track of a matrix $\Gamma$ initialized with $\Gamma = \lambda \mathbf{I}$, $\mathbf{I} \in \mathbb{R}^{N \times N}$ is an identity matrix, where $\lambda > 0$ is a regularization parameter and $N$ is the parameter size of the neural network. Assuming a Gaussian distributed reward posterior, the variance of the distribution is estimated by
\begin{equation}
    \sigma^2_{t, A} = \lambda g_{\theta}^T(x_{t, A})\Gamma^{-1}g_{\theta}^T(x_{t, A}),
\end{equation}
where $g_\theta$ is the gradient of the output with respect to $\theta$ for the entire neural network. The agent then chooses action
\begin{equation}
    A_t = \arg\max_{a \in \mathcal{A}}\hat{R}_{t, a}, \qquad \hat{R}_{t, a} \sim \mathcal{N}(f_{\theta}(x_{t, a}), \alpha \sigma^2_{t, a}),
\end{equation}
where $\alpha$ is an exploration hyperparameter and $f_{\theta}$ is the neural network. The agent updates $\Gamma \leftarrow \Gamma + g_{\theta}(x_{t, A})g_{\theta}(x_{t, A})^T / m$ with $m$ being the width of the neural network, assuming all layers with the same width. To the best of the authors' knowledge, the neural Thompson sampling strategy is not well adopted by industry due to its heavy computational complexity. 

\subsection{Upper Confidence Bound (UCB), Its Extensions and Related Adoptions}
Upper confidence bound (UCB) is a general optimism facing uncertainty exploration strategy where the agent tends to choose actions that it has more uncertainties about. The purpose of doing so is to gather information that could help either identify the best action or eliminate bad actions. Similar to the subsection above, we first consider a vanilla UCB algorithm \cite{agrawal1995sample, auer2002finite} for multi-armed bandits. A vanilla UCB agent selects an action $A_t$ by 
\begin{equation}
    A_t = \arg\max_{a \in \mathcal{A}} \left(\frac{\sum_{\tau=1}^{t} R_\tau \mathbbm{1}(A_\tau = a)}{n_{t, a}} + \alpha\sqrt{\frac{\ln(t)}{n_{t, a}}}\right),
\end{equation}
where $n_{t, A} = \sum_{\tau = 1}^t \mathbbm{1}(A_\tau = A)$ and $\alpha \in \mathbb{R_+}$ is a hyperparameter. There have been multiple lines of work analyzing vanilla UCB's impact on production recommender systems in terms of degenerating feedback loops \cite{jiang2019degenerate, Guo2023evaluate}, modeling attritions \cite{ben2022modeling} and its extension to Collaborative Filtering \cite{nakamura2015ucb}.

The immediate extension to UCB to linear bandit problems is LinUCB \cite{li2010contextual, chu2011contextual}. A LinUCB agent initializes two variables for any new action $A$ that it has never seen before, $\Gamma_A = \mathbf{I}_d$ ($d$-dimensional identity matrix) and $b_A = \mathbf{0}_{d\times 1}$ ($d$-dimensional zero vector), where $d$ is the linear parameter size. At time step t, given a new context representation $\psi_{S_t}$ and each action $\phi_{A}$, the agent concatenates the two vectors to get a context-action pair representation $x_{t, A} = \text{concat}(\psi_{S_t}, \phi_A)$. The agent selects action
\begin{equation}
    A_t = \arg\max_{a \in \mathcal{A}}\left((\Gamma^{-1}_a b_a)^T x_{t, a} + \alpha \sqrt{x_{t, a}^T\Gamma_a^{-1}x_{t, a}}\right).
\end{equation}
After receiving $R_{t+1}$, the agent then updates $\Gamma_{A_t}$ and $b_{A_t}$ with $\Gamma_{A_t} \leftarrow \Gamma_{A_t} + x_{t, A_t}x_{t, A_t}^T$ and $b_{A_t} \leftarrow b_{A_t} + R_{t+1} x_{t, A_t}$. LinUCB is a popular algorithm for adoption in many industry usecases and particularly in the space of recommender systems \cite{li2010contextual, zhao2013interactive, qin2014contextual, wang2017biucb}. 

With the rise of neural networks and deep learning, so comes the need for optimism based exploration neural methods. NeuralUCB \cite{zhou2020neural} and NeuralLinUCB \cite{xu2021neural} are two representing examples of such algorithms. We first introduce NeuralLinUCB as it's a natural extension of LinUCB. Given a neural network representation, NeuralLinUCB no longer initiates a matrix $\Gamma_A$ for each action $A$ and instead keeps track of a single matrix $\Gamma$ initialized with $\lambda \mathbf{I}$ where $\lambda > 0$ is a regularization factor and $\mathbf{I} \in \mathbb{R}^{N \times N}$ is an identity matrix. Suppose the neural network's last layer representation is $\sigma_{\theta}(x_{t, A}) \in \mathbb{R}^N$ and full model output $f_\theta(x_{t, A}) \in \mathbb{R}$, the agent takes action
\begin{equation}
    A_t = \arg\max_{a \in \mathcal{A}}\left(f_{\theta}(x_{t, a}) + \alpha \sqrt{\sigma_{\theta}(x_{t, a})^T\Gamma^{-1}\sigma_{\theta}(x_{t, a})}\right),
\end{equation}
with the same update for $\Gamma$ and $b$ except that now the update happens with the representation $\sigma_\theta(x_{t,A})$ instead of $x_{t,A}$. NeuralUCB follows the same principle as above, but replacing $\sigma_{\theta}(x_{t,A})$ with the gradients of all parameters within the neural network. 

Although empirically the above neural methods, Neural Thompson sampling, Neural UCB and Neural LinUCB, have achieved good performance on synthetic datasets, their downside is that they all require inverting a square matrix with its dimension equal to either the entire neural network's parameter size or the size of the last layer representation. Both in many cases are intractable, especially for real-world environments and complex neural networks. 